% README for PAGINA-WED (LaTeX format)
\documentclass[a4paper,12pt]{article}
\usepackage[spanish]{babel}
\usepackage[utf8]{inputenc}
\usepackage{hyperref}
\title{Documentaci\'on del Proyecto \texttt{PAGINA-WED}}
\author{Desarrollador}
\date{\today}
\begin{document}
\maketitle
\section{Visi\'on general}
El proyecto \texttt{PAGINA-WED} es una aplicaci\'on web de comercio electr\'onico escrita principalmente en PHP, que se ejecuta sobre un servidor Apache con MySQL como sistema gestor de bases de datos (entorno XAMPP o similar).
Ofrece un conjunto completo de funcionalidades: administraci\'on de usuarios, gesti\'on de productos, carrito de compras, registro/login, perfiles, mensajer\'ia interna y panel de control para administradores.
Para facilitar el mantenimiento y la evoluci\'on, la l\'ogica est\'a organizada en carpetas que separan claramente la capa de cliente, el backend de administraci\'on, los endpoints API y los m\'odulos de configuraci\'on y conexi\'on.

\section{Estructura de carpetas}
\begin{itemize}
  \item \texttt{admin/}: Interfaces y scripts para administradores, CRUD de usuarios y productos, panel.
  \item \texttt{api/}: Endpoints utilizados por el cliente y la aplicaci\'on para operaciones as\'incronas (registro, login, actualizaci\'on, etc.).
  \item \texttt{client/}: P\'aginas de usuario finales (carrito, perfil, hist\'orial, compras).
  \item \texttt{core/}: Funciones centrales como conexi\'on a base de datos, configuraci\'on, manejo de sesiones y cierre de sesi\'on.
  \item \texttt{database/}: Archivo SQL con la definici\'on de la base de datos \texttt{negocio\_web.sql}.
  \item \texttt{pages/}: P\'aginas p\'ublicas (contacto, login, crear cuenta, ofertas).
  \item \texttt{img/}, \texttt{vendor/}: Recusos est\'aticos e incluye la librer\'ia PHPMailer.
\end{itemize}

\section{Tecnolog\'ias y metodolog\'ias}
Desde una perspectiva de arquitecturas senior, se emplea una base supletoria de patrones y librer\'ias para maximizar mantenibilidad y seguridad.
\begin{itemize}
  \item El c\'odigo est\'a escrito en PHP en estilo procedural con un enfoque de inclusiones y archivos de configuraci\'on que ejercen una funci\'on similar a la capa de "modelo". Aunque no se sigue un framework MVC completo, la separaci\'on de responsabilidades est\'a clara en carpetas y archivos dedicados.
  \item Se utiliza la extensi\'on nativa \texttt{mysqli} para interactuar con la base de datos MySQL; todas las llamadas a la base de datos pasan por la funci\'on \texttt{conexion()} definida en \texttt{core/conexion.php}. Esto proporciona encapsulamiento, minimiza la duplicaci\'on de c\'odigo y facilita futuros cambios de motor de datos.
  \item La l\'ogica de sesiones y autenticaci\'on se encapsula dentro de \texttt{core/sesiones.php}. Dicho archivo inicializa session\_start(), verifica inyecciones de identificador y expone helpers como \texttt{isLogged()} o \texttt{esAdmin()} para controlar acceso. El uso de esta capa oculta detalles de implementaci\'on y permite ajustar par\'ametros (timeout, regeneraci\'on de ID) en un solo lugar.
  \item Para estilizado y componentes en el frontend se aprovecha Tailwind CSS a trav\'es de su CDN (ver referencias en las p\'aginas \texttt{index1.php}, \texttt{client/*} y \texttt{core/configuracion.php}). Este enfoque evita dependencias locales y simplifica el dise\~no responsivo.
  \item Se integran librer\'ias de terceros cuando son necesarias: en particular PHPMailer (ubicada en \texttt{vendor/PHPMailer/}) se utiliza para env\'io de correos desde \texttt{core/smtp\_config.php}. El uso de paquetes autogenerados mediante composer puede ser implementado en futuras iteraciones.
  \item Las APIs internas siguen el estilo RESTful b\'asico; reciben peticiones POST/GET, gestiona encabezados y devuelven respuestas JSON con estados estandarizados. Cada endpoint reside en la carpeta \texttt{api/} y reusa la conexi\'on y validaci\'on de sesiones.
  \item Se aplican buenas pr\'acticas de seguridad: sanitizaci\'on de entradas (\texttt{filter\'_input}, \texttt{mysqli\'_real\'_escape\'_string}), contrase\~nas hasheadas con \texttt{password\'_hash()}, y revisiones en inclusiones para evitar direcciones relativas manipulables.
  \item El c\'odigo HTML embebido y los scripts JavaScript (peque
tilas validaciones, dynamic injection) se mantienen dentro de los mismos archivos de vista para proyectos de bajo volumen; en un futuro, se podr\'ia extraerlos a archivos est\'aticos o usar un transpiler.
\end{itemize}

\section{Librer\'ias y dependencias}
El proyecto utiliza algunos recursos externos que se gestionan manualmente o mediante distribuci\'on CDN. Las m\'as importantes son:
\begin{itemize}
  \item \textbf{Tailwind CSS}: se carga directamente desde el CDN oficial en las p\'aginas, lo que permite aprovechar su sistema de utilidades para rapidez de desarrollo sin instalarla localmente.
  \item \textbf{PHPMailer}: la versi\'on incluida en \texttt{vendor/PHPMailer/} se utiliza para manejo de correo saliente. La clase se carga con \texttt{require\_once} en \texttt{core/smtp\_config.php}. La configuraci\'on SMTP (host, puerto, credenciales) se parametriza en el mismo archivo.
  \item \textbf{Sin Composer}: actualmente no se emplea un gestor de paquetes PHP, por lo que las dependencias se agregan a mano en el directorio \texttt{vendor/}. Esto simplifica el despliegue pero deja una mejora pendiente.
  \item \textbf{Recursos JavaScript}: no existen bibliotecas JS externas significativas; cualquier comportamiento del lado del cliente se implementa mediante snippets embebidos.
  \item \textbf{Base de datos y SQL}: las definiciones de esquema se encuentran en \texttt{database/negocio\_web.sql}. No se utiliza ORM, lo que hace transparente el acceso pero exige rigor en la preparación de consultas.
\end{itemize}

\section{Funcionalidades destacadas}
\begin{enumerate}
  \item Registro y validaci\'on de usuarios.
  \item Inicio de sesi\'on con contra\'se\~na cifrada y control de sesiones.
  \item Panel de administraci\'on para crear, editar y eliminar usuarios y productos.
  \item Carrito de compras con historial y confirmaci\'on de compra.
  \item P\'aginas de perfiles y mensajer\'ia entre usuarios.
  \item APIs para operaciones cr\'iticas (actualizar perfil, cambiar contrase\~na, eliminar cuenta, etc.).
\end{enumerate}

\section{Buenas pr\'acticas}
\begin{itemize}
  \item Separaci\'on de responsabilidades en carpetas.
  \item Uso de inclusiones para encapsular la conexi\'on a la base de datos y configuraciones.
  \item Sanitizaci\'on de entradas en los scripts API y formularios del lado del servidor.
  \item Estructura de la base de datos normalizada para usuarios, productos, pedidos y mensajes.
  \item Comentarios en c\'odigo y nomenclatura clara para variables y funciones; se prefieren nombres en espa\~nol en el servidor por consistencia.
  \item Revisa\'on manual de rutas relativas en inclusiones para prevenir archivos accesibles directamente desde la web.
  \item Manejo de errores simple pero efectivo: se retornan mensajes de error JSON en APIs y se registran excepciones de la base de datos en un archivo de log (no incluido).
  \item Evitar duplicaci\'on de HTML mediante fragmentos de cabecera/pie reutilizables (no desarrollados completamente, pero la idea est\'a presente).
\end{itemize}

\section{Instalaci\'on}
\begin{enumerate}
  \item Colocar los archivos en un servidor con PHP y MySQL (ej. XAMPP bajo \\htdocs).
  \item Crear la base de datos importando \texttt{database/negocio\_web.sql} mediante phpMyAdmin.
  \item Configurar los datos de conexi\'on en \texttt{core/conexion.php} y ajustes adicionales en \texttt{core/configuracion.php}.
  \item Asegurar permisos de escritura en carpetas necesarias (si aplica).
\end{enumerate}

\section{Nota final}
Este README en LaTeX ofrece una visi\'on global y se puede compilar a PDF usando \texttt{pdflatex README.tex}.
Se recomienda revisar y mejorar la seguridad antes de desplegar en producci\'on.

\end{document}
